\documentclass[10pt,a4paper]{amsart}
\usepackage[margin=19mm]{geometry}
\usepackage{amsmath,amssymb,mathtools}
\usepackage[scr=rsfs,cal=euler]{mathalfa}
\usepackage{xfrac}
\usepackage[unicode,hidelinks,bookmarks=false]{hyperref}
\newcommand{\ie}{i.e.\ }
\newcommand{\cf}{cf.\ }
\newcommand{\resp}{respectively\ }
\newcommand{\Subsection}{\subsection}
\providecommand{\nobreakdash}{\nobreak\hbox{-}}
\setlength{\emergencystretch}{2em}
\multlinegap0pt
\usepackage{bm}
\usepackage{bbm}
\usepackage{dsfont}
\usepackage{xspace}
\usepackage{cancel}
\usepackage{mathtools}
\usepackage{amsthm}
\makeatletter
\@ifundefined{theorem}{\newtheorem{theorem}{Theorem}}{}
\@ifundefined{lemma}{\newtheorem{lemma}[theorem]{Lemma}}{}
\makeatother
\title[Labeled Chip-Firing on Directed $k$-ary Trees and Where Chip]{Labeled Chip-Firing on Directed $k$-ary Trees and Where Chips Land}
\author{Ryota Inagaki, Tanya Khovanova, Austin Luo}
\date{Source: arXiv:2506.20656v1 (CC BY 4.0)}
\begin{document}
\maketitle

\noindent\emph{Source and licence.} Labeled Chip-Firing on Directed $k$-ary Trees and Where Chips Land. Version arXiv:2506.20656v1 (CC BY 4.0), distributed under the Creative Commons Attribution 4.0 International licence, as recorded in the arXiv licence metadata for this exact version and shown on the abstract page https://arxiv.org/abs/2506.20656v1. Full rights notice: licenses/arxiv-2506.20656-CC-BY-4.0.txt.\par\smallskip

\noindent\textit{Editorial scope.} Counted unit: the Theorem whose source label is \texttt{thm:ranges} in the article (source0.tex lines 452--456, source label \texttt{thm:ranges}), delivered together with the article's own complete original proof of it (source0.tex lines 458--493, one \texttt{proof} environment of 36 lines). No mathematics is written, repaired or re-derived: statement and proof are the article's own text line for line. Required context retained uncounted: lem:Ranges (lines 434--436). Every definition, display and label of those lines is retained unchanged. Omitted from this package and not redistributed: 1--433, 437--451, 457--457, 494--950. Nothing inside a retained excerpt is omitted or altered: every retained body is delivered exactly as the article's lines are written, so no omission receipt of any kind accompanies this package. Citations: the retained excerpts carry no citation command at all, so no bibliography entry of the article is transcribed and no reference is redistributed. Reference adaptation: none was needed and none was made; every reference inside the delivered unit resolves inside it, so no unresolved reference or citation is produced. Printed slips kept literal: no printed word, symbol, hypothesis or number of the counted statement or of its proof was altered. Layout and class: the article's own class is \texttt{article} with packages including amsmath, amssymb, amscd, amsthm, amsfonts, subfig, graphicx, hyperref, verbatim, float, mathtools, color, colortbl, caption; the wrapper is the lane's pinned \texttt{amsart} editorial wrapper at 10pt on A4, which restates the theorem-like environments the retained text needs and adds the article's own macro definitions that the retained text needs (none), with extra wrapper packages (bm, bbm, dsfont, xspace, cancel, mathtools). The delivered numbering is the unit's own. Assets: no retained span contains a figure environment, a graphics-inclusion command, a PDF-page inclusion, a table or a TikZ picture; no image file is copied into this package, the package declares no asset dependency and no image file is redistributed with it. Licence: version arXiv:2506.20656v1 of this article is distributed under the Creative Commons Attribution 4.0 International licence, as recorded in the arXiv licence metadata for this exact version and shown on the abstract page https://arxiv.org/abs/2506.20656v1. A full rights notice is delivered at licenses/arxiv-2506.20656-CC-BY-4.0.txt. Characters: every retained excerpt is pure ASCII, so the delivered engine, XeLaTeX with its default Latin Modern text font, reports no missing character.
\begin{lemma}\label{lem:Ranges}
   Given a landing $(t,x)$ on layer 2 and chip $c$ such that $a(t,x) \leq c \leq b(t,x)$, there exists a firing strategy that places $c$ at $(t,x)$.
\end{lemma}

\begin{theorem}
\label{thm:ranges}
    For any landing $(t,x)$ on any layer, any chip in the range from $a(t,x)$ to $b(t,x)$ can appear there. In particular, the set of chips that can land at a particular vertex $v_t$ on the last layer forms a consecutive landing range
    \[a(t, 1) = \prod_{j=1}^n t_j \quad \textrm{ to } \quad b(t, 1) = k^n+1-\prod_{j=1}^n(k+1-t_j)\] inclusive. 
\end{theorem}

\begin{proof}
    We already proved the base case for layer 2 in Lemma~\ref{lem:Ranges}. We proceed by induction.
    
    Assume that at layer $i+1$, the chips landing at landing order $(t,x)$ form the interval $[a(t,x), b(t,x)]$.

    We now prove that the same is true for layer $i+2$. Consider a landing order $(t',x')$ at layer $i+2$, where $t'=t_1t$, where $t$ is a traversal string of length $i$. We now consider a subtree starting at a child $v_{t_1}$ of the root. Relative to this subtree, our landing order is on layer $i+1$. Thus, by the induction hypothesis, our chip can only arrive from the interval of orders $(t',(a(t,x))$ to $(t',(b(t,x))$.
    
    What is left to prove is that chips that can land at this interval form an interval. Consider two consecutive landing orders $(t_1,x)$ and $(t_1,x+1)$ on layer 2. Assume that $n > 1$ and $x < k^{n-1}$, since otherwise $x+1$ would be greater than the number of chips in each vertex in layer 2 and $(t_1, x+1)$ would not be a valid landing order.  We know that 
    \[a(t_1,x) = xt_1 < a(t_1,x+1) = (x+1)t_1\]
    \[b(t_1,x) = k^n +1- (k^{n-1} + 1 - x) (k+1-t_1) < b(t_1,x+1) =k^n+1 - (k^{n-1} + 1 - x-1) (k+1-t_1).\]
    Now, we want to show that for $x \in \{1, 2, \dots, k^{n-1}-1\}$,
    \[a(t_1,x+1) \leq b(t_1,x) + 1,\]
    or equivalently,
     \begin{equation}\label{eq:SlightlyBoringInequality}
         (x+1)t_1 \leq k^n +1- (k^{n-1} + 1 - x) (k+1-t_1) + 1.
     \end{equation}
     First, we rearrange the right-hand side
     \begin{multline*}
        k^n +1- (k^{n-1} + 1 - x) (k+1-t_1) + 1 = k^n-k^n + k^{n-1} (t_1-1) -k-1+t_1 +x (k+1-t_1) + 2 = \\
        k^{n-1} (t_1-1) -k+ t_1 +x (k+1-t_1)+1,
     \end{multline*}
     implying that inequality (\ref{eq:SlightlyBoringInequality}) above is equivalent to
     \[(x+1)t_1 \leq k^{n-1} (t_1-1) -k+ t_1 +x (k+1-t_1)+1.\]
     Rearranging, we obtain the following
     \begin{equation}\label{eq:SimplerInequality}
         (2t_1-k-1) x+k-1 \leq (t_1-1) k^{n-1}.
     \end{equation}

     Suppose $t_1 \leq k-1$, then we have
     \[(2t_1-k-1) x+k-1 \leq (t_1-2) x+k-1 \leq (t_1-2)(k^{n-1} -1) + k-1).\]
    The right-hand side can be rearranged to
    \[(t_1 - 1)k^{n-1} - k^{n-1} + k  \leq  (t_1 - 1)k^{n-1}.\]
    Suppose $t_1 = k$, then we have 
     \[(2t_1-k-1) x+k-1 = (k-1) x+k-1 = (k-1)(x+1) \leq (k-1) k^{n-1} = (t_1-1) k^{n-1}.\]
     We showed that two consecutive landing orders at the same vertex overlap. It follows that the union of any number of consecutive landing orders at the same vertex form an interval.
\end{proof}

\end{document}
